% !TeX root = ../thuthesis-example.tex

\chapter{系统设计与实现（实现步骤截图、环境安装、代码实现）}

\section{总体架构}

系统采用“数据层、计算层、服务层、展示层”四层结构。数据层包括本地 Kaggle 原始 CSV、本地 1\%/5\% 样本、HDFS CSV 输入和 HDFS Parquet 清洗层；计算层包括 Spark SQL pipeline、AQE、推荐/关联/异常等算法模块和 Spark History event log；服务层包括 Flask API、SQLite 作业记录、run manifest 与 benchmark evidence 服务；展示层包括 React 前端业务页面、质量页和 Ops 页。

图\ref{fig:yarn-architecture}展示了 Docker YARN 实验集群拓扑。该拓扑保留了 Hadoop 课程设计所需的 NameNode、DataNode、ResourceManager、NodeManager 和 Spark History Server，同时增加 spark-client 作为提交入口。前端与后端按 \texttt{start-dev.sh} 启动：Flask 后端为 \texttt{5051}，Vite 前端为 \texttt{5173}；HDFS UI、YARN UI 和 Spark History UI 分别映射到 \texttt{19870}、\texttt{18088} 和 \texttt{28080}，避免与已有 \texttt{master} 容器冲突。

\begin{figure}[htbp]
    \centering
    \includegraphics[width=0.95\textwidth]{yarn_lab_architecture.png}
    \caption{Docker ARM64 Hadoop YARN 实验集群拓扑（由 chat-gpt-image-2 生成）}
    \label{fig:yarn-architecture}
\end{figure}

\section{环境安装与启动步骤}

\subsection{ARM64 Hadoop 镜像构建}

本项目避免使用 amd64 Hadoop 镜像，也避免把整个 \texttt{master} 容器 commit 成大镜像。最终方案是自建精简 ARM64 Hadoop 镜像，只包含运行 Hadoop 所需的 Java、Hadoop 文件、配置和启动脚本。核心思路如下：

\begin{lstlisting}[style=codexcode]
FROM ubuntu:24.04

RUN apt-get update && apt-get install -y \
    openjdk-8-jdk-headless bash ssh rsync procps curl \
    && rm -rf /var/lib/apt/lists/*

COPY hadoop /opt/hadoop

ENV HADOOP_HOME=/opt/hadoop
ENV HADOOP_CONF_DIR=/opt/hadoop/etc/hadoop
ENV PATH=/opt/hadoop/bin:/opt/hadoop/sbin:$PATH
\end{lstlisting}

与 \texttt{docker commit master} 相比，该方法不会复制 HDFS 数据目录、日志、Hive、HBase 和临时文件，因此镜像体积更可控，也符合 ARM Mac 的架构约束。

\subsection{Docker Compose 启动}

集群通过 \texttt{yarn-lab} profile 启动。常用命令如下：

\begin{lstlisting}[style=codexcode]
docker compose --profile yarn-lab up -d
docker compose --profile yarn-lab ps
docker compose --profile yarn-lab config --quiet
OPEN_BROWSER=0 ./start-dev.sh
\end{lstlisting}

端口规划如表\ref{tab:ports}所示。Hadoop UI 端口均为偏移端口，原因是用户原有 \texttt{master} 容器已经占用 Hadoop 默认 UI 端口；前后端访问地址则以 \texttt{start-dev.sh} 为准。

\begin{table}[htbp]
\centering
\caption{实验环境端口规划}
\label{tab:ports}
\begin{tabularx}{0.92\textwidth}{l l X}
\toprule
\textbf{服务} & \textbf{宿主机端口} & \textbf{说明} \\
\midrule
Flask API & \texttt{5051} & \texttt{start-dev.sh} 启动的后端接口，前端通过该地址访问 \texttt{/api/v1}。 \\
Vite 前端 & \texttt{5173} & \texttt{start-dev.sh} 启动的正确前端访问地址，用于打开 \texttt{/ops}。 \\
NameNode UI & \texttt{19870} & 查看 HDFS 文件、block 和 DataNode 状态。 \\
ResourceManager UI & \texttt{18088} & 查看 YARN application、队列和运行状态。 \\
Spark History Server & \texttt{28080} & 查看 Spark event log 和 stage/job 信息。 \\
\bottomrule
\end{tabularx}
\end{table}

\subsection{HDFS 目录准备}

启动后需要创建实验目录，并上传样本数据。正式样本来自真实 Oct/Nov 原始 CSV 的用户级采样，不再使用 87.4KB smoke 文件作为论文结论依据。

\begin{lstlisting}[style=codexcode]
docker exec enormous-data-yarn-namenode hdfs dfs -mkdir -p /user/course
docker exec enormous-data-yarn-namenode hdfs dfs -mkdir -p /spark-history
docker exec enormous-data-yarn-namenode hdfs dfs -put -f \
  /data/sample/ecommerce_user_sample_1pct.csv \
  /user/course/ecommerce_behavior_user_sample_1pct/
docker exec enormous-data-yarn-namenode hdfs dfs -put -f \
  /data/sample/ecommerce_user_sample_5pct.csv \
  /user/course/ecommerce_behavior_user_sample_5pct/
\end{lstlisting}

清理后的 HDFS 关键路径包括 1\%/5\% CSV 输入、YARN-only/AQE/algorithm/Parquet 的处理输出，以及 6 个最终 Spark History event log。旧的 \texttt{application\_178098*} 和 \texttt{application\_1780991452919\_0001} 到 \texttt{0005} 大日志已经清理，避免继续触发 Spark 3.5.5 History Server 的 JSON 长字符串限制。

\section{Spark 提交链路实现}

\subsection{YARN client 配置}

第一阶段采用 YARN client mode。这样 driver 运行在 spark-client 容器，Flask 仍可读取本地 \texttt{data/cache}，前端也能继续使用原有缓存文件。核心配置包括：

\begin{lstlisting}[style=codexcode]
spark:
  master: yarn
  deployMode: client
  appName: ecommerce-yarn-lab
  config:
    spark.sql.adaptive.enabled: true
    spark.sql.adaptive.skewJoin.enabled: true
    spark.sql.shuffle.partitions: 64
    spark.sql.maxPlanStringLength: 8192
    spark.eventLog.enabled: true
    spark.eventLog.dir: hdfs:///spark-history
\end{lstlisting}

\texttt{spark.sql.maxPlanStringLength} 是本次修复的重要配置\cite{spark_configuration}。旧 5\% event log 在 Spark History Server 的 \texttt{/stages} API 上会触发 JSON 长字符串限制，限制 plan string 后重新运行的最终 event log 体积明显降低，便于长期保留。

\subsection{提交脚本}

提交前先将 \texttt{spark\_jobs} 打包为 zip，再通过 \texttt{spark-submit} 提交到 YARN：

\begin{lstlisting}[style=codexcode]
python scripts/package_spark_jobs.py
spark-submit \
  --master yarn \
  --deploy-mode client \
  --py-files build/spark_jobs.zip \
  spark_jobs/main.py \
  --config configs/yarn-client.yaml
\end{lstlisting}

Flask 的 refresh 不再直接执行 \texttt{python -m spark\_jobs.main}，而是调用 submit runner。运行完成后，\texttt{SparkPipelineRunner} 从 stdout 中解析 \texttt{Spark run manifest} 路径，读取 manifest 并写入 SQLite 作业表。为了避免后端消息过长，runner 只保留 stdout/stderr tail，完整日志写入 \texttt{data/cache/run\_logs}。

\subsection{YARN cluster mode 入口}

在第一阶段，系统默认使用 YARN client mode，因为 Flask 与 React 仍读取 spark-client 容器中的 \texttt{data/cache}。同时项目新增了 \texttt{configs/yarn-cluster.yaml} 与 \texttt{scripts/submit\_yarn\_cluster.sh}，作为第二阶段 cluster mode 入口。cluster mode 会把 driver 交给 NodeManager 运行，并通过 \texttt{--files} 分发配置文件，通过 \texttt{spark.yarn.appMasterEnv.PYSPARK\_PYTHON} 和 \texttt{spark.executorEnv.PYSPARK\_PYTHON} 固定 Python 3.11，避免 driver 与 executor Python 小版本不一致。由于 cluster mode 下本地 JSON cache 不再天然落在 Flask 可读目录，后续若要把 cluster mode 设为默认刷新路径，需要继续实现 HDFS 输出读取或缓存同步。

\section{核心算法与内存优化实现}

\subsection{禁止 full collect}

早期推荐模块会把 \texttt{recommendation\_features.collect()} 的完整结果带回 driver，再生成大 JSON。这在小样本下能运行，但数据量扩大后会直接增加 driver OOM 风险。优化后，summary、coverage、fallback、duplicate 和 confidence 等质量指标全部使用 Spark 聚合完成；前端展示只收集 preview limit 行数；完整明细写入 Parquet。

\begin{lstlisting}[style=codexcode]
preview_rows = (
    recommendation_features
    .orderBy(F.desc("score"))
    .limit(preview_limit)
    .collect()
)

quality = recommendation_features.agg(
    F.count("*").alias("candidate_rows"),
    F.avg("confidence").alias("avg_confidence"),
    F.sum(F.col("is_fallback").cast("int")).alias("fallback_rows"),
)
\end{lstlisting}

\subsection{WindowExec 警告优化}

Spark 日志中曾出现大量 \texttt{WindowExec: No Partition Defined} 警告。这类警告表示窗口函数没有分区键，Spark 会把数据移动到单分区处理，小样本时只是警告，全量时会变成瓶颈。优化后，ranking 逻辑统一通过 helper 添加全局 rank 或业务分区 rank，避免模块中散落无分区 \texttt{Window.orderBy()}。

\subsection{关联分析组合膨胀控制}

商品关联图谱最容易产生二次级组合膨胀。若某个 session 中商品数过多，直接两两配对会产生 $O(n^2)$ 行。系统加入 \texttt{top\_nodes}、\texttt{max\_pair\_products\_per\_session} 和 pair explosion 指标。质量报告记录 \texttt{eligible\_session\_count}、\texttt{pair\_base\_rows} 和 \texttt{pair\_rows\_per\_input\_row}，用于判断数据规模增长是否仍受控。

\subsection{Parquet 清洗层}

CSV 原始层只作为输入层。清洗后的事件表写入 HDFS Parquet，并按日期字段分区。Parquet 输入对照组直接读取清洗目录，跳过 CSV schema 解析和清洗阶段的部分开销。图\ref{fig:pipeline-flow}展示了从原始 CSV 到前端 Ops 证据的处理流程。

\begin{figure}[htbp]
    \centering
    \includegraphics[width=0.95\textwidth]{pipeline_optimization_flow.png}
    \caption{Spark 数据处理与优化流程（由 chat-gpt-image-2 生成）}
    \label{fig:pipeline-flow}
\end{figure}

\section{前端 Ops 页实现}

前端新增 \texttt{/api/v1/ops/evidence} 接口调用。后端 \texttt{BenchmarkEvidenceService} 读取正式 benchmark 目录，只返回精简字段：样本、对照组、耗时、吞吐、YARN app id、质量状态、driver 峰值内存和结果路径。前端 Ops 页新增两块内容：第一块是 benchmark 表，用于展示 1\% 五组和 5\% 两组实验；第二块是 HDFS 输入与数据清理策略，用于证明正式输入、Parquet 层和保留的 Spark History app id。

该设计解决了两个答辩痛点。第一，用户不需要在终端手动翻 \texttt{data/benchmarks}；第二，前端不会读取 Spark stderr 长文本，因此不会把 History Server JSON 限制问题转移到浏览器。核心接口返回结构如下：

\begin{lstlisting}[style=codexcode]
{
  "benchmark_runs": [
    {
      "sample": "1pct",
      "variant": "yarn_only_csv",
      "elapsed_seconds": 964.889,
      "rows_per_second": 1143.869,
      "spark_application_id": "application_1780991452919_0016",
      "quality_status": "passed"
    }
  ],
  "hdfs_inputs": [...],
  "cleanup_policy": {...}
}
\end{lstlisting}

图\ref{fig:ops-evidence-panel}展示了前端 Ops 页的证据集成方式。后端将 benchmark 结果、HDFS 输入快照、Spark History app id 和清理策略统一封装成轻量 JSON，前端只负责读取和展示这些证据，不直接解析 Spark 长日志。

\begin{figure}[htbp]
    \centering
    \includegraphics[width=0.95\textwidth]{ops_evidence_panel.png}
    \caption{前端 Ops 证据面板集成示意图（由 chat-gpt-image-2 生成）}
    \label{fig:ops-evidence-panel}
\end{figure}

\section{旧数据清理实现}

清理工作遵循“只删生成物，不删原始数据”的原则。保留的数据包括：

\begin{itemize}
    \item \texttt{data/raw/kaggle/ecommerce\_behavior/2019-Oct.csv}，约 5.3GB。
    \item \texttt{data/raw/kaggle/ecommerce\_behavior/2019-Nov.csv}，约 8.4GB。
    \item \texttt{data/sample/ecommerce\_user\_sample\_1pct.csv}，约 141.5MB。
    \item \texttt{data/sample/ecommerce\_user\_sample\_5pct.csv}，约 694.3MB。
    \item 三个正式 benchmark 目录：1\% fixed baseline、1\% YARN matrix、5\% YARN matrix。
    \item 最终 Spark History event log：\texttt{application\_1780991452919\_0016} 到 \texttt{0021}。
\end{itemize}

删除的数据包括早期 smoke、retry、history recheck benchmark 目录，旧 \texttt{/user/course/ecommerce\_behavior\_processed} 小输出，以及旧的大 Spark History 日志。清理后本地 benchmark 目录收敛到约 292KB，HDFS \texttt{/spark-history} 只保留 6 个最终 event log，约 87MB。存储快照如图\ref{fig:storage-cleanup}所示。

\begin{figure}[htbp]
    \centering
    \includegraphics[width=0.9\textwidth]{storage_cleanup_snapshot.png}
    \caption{旧数据清理后的存储治理示意图（由 chat-gpt-image-2 生成，精确容量见正文说明）}
    \label{fig:storage-cleanup}
\end{figure}

\section{实验结果}

\subsection{Benchmark 对照表}

图\ref{fig:experiment-matrix}展示了本课题的实验对照矩阵。正式实验结果如表\ref{tab:benchmark-results}所示。1\% 样本包含五组对照，5\% 样本包含算法优化 CSV 与 Parquet 两组压力验证。所有 YARN application 均为 \texttt{SUCCEEDED}。

\begin{figure}[htbp]
    \centering
    \includegraphics[width=0.95\textwidth]{experiment_matrix_methodology.png}
    \caption{YARN + Spark 实验对照矩阵（由 chat-gpt-image-2 生成，精确数值见表\ref{tab:benchmark-results}）}
    \label{fig:experiment-matrix}
\end{figure}

\begin{table}[htbp]
\centering
\caption{1\% 与 5\% benchmark 实验结果}
\label{tab:benchmark-results}
\scriptsize
\begin{tabularx}{0.98\textwidth}{l l r r r X}
\toprule
\textbf{样本} & \textbf{对照组} & \textbf{耗时(s)} & \textbf{输入行} & \textbf{输出行} & \textbf{Application ID} \\
\midrule
1\% & Baseline local CSV & 198.843 & 1,104,709 & 1,103,488 & \texttt{local-1781015058333} \\
1\% & YARN-only CSV & 964.889 & 1,104,709 & 1,103,488 & \texttt{application\_1780991452919\_0016} \\
1\% & YARN + AQE CSV & 538.657 & 1,104,709 & 1,103,488 & \texttt{application\_1780991452919\_0017} \\
1\% & YARN + algorithm CSV & 535.804 & 1,104,709 & 1,103,488 & \texttt{application\_1780991452919\_0018} \\
1\% & YARN + Parquet & 573.183 & 1,103,488 & 1,103,488 & \texttt{application\_1780991452919\_0019} \\
5\% & YARN + algorithm CSV & 1875.099 & 5,412,207 & 5,405,748 & \texttt{application\_1780991452919\_0020} \\
5\% & YARN + Parquet & 1828.329 & 5,405,748 & 5,405,748 & \texttt{application\_1780991452919\_0021} \\
\bottomrule
\end{tabularx}
\end{table}

\subsection{Spark History 指标}

旧 Spark History Server 的 \texttt{/stages} API 在部分应用上会返回 500。为避免指标缺失，系统新增 \texttt{scripts/collect\_eventlog\_metrics.py}，直接解析 HDFS 中的 rolling event log \texttt{.zstd} 文件，汇总 task、shuffle 和 spill 指标。表\ref{tab:history-results}给出 6 个正式 YARN 应用的补采集结果。

\begin{table}[htbp]
\centering
\caption{Spark event log 指标补采集结果}
\label{tab:history-results}
\scriptsize
\begin{tabularx}{0.98\textwidth}{l r r r r r}
\toprule
\textbf{Application ID} & \textbf{Tasks} & \textbf{Failed} & \textbf{Retried} & \textbf{Shuffle Read} & \textbf{Memory Spill} \\
\midrule
\texttt{application\_1780991452919\_0016} & 6,290 & 0 & 0 & 9.4GB & 2.1GB \\
\texttt{application\_1780991452919\_0017} & 4,526 & 0 & 0 & 9.4GB & 728.0MB \\
\texttt{application\_1780991452919\_0018} & 4,526 & 0 & 0 & 9.4GB & 728.0MB \\
\texttt{application\_1780991452919\_0019} & 5,057 & 0 & 0 & 9.4GB & 792.5MB \\
\texttt{application\_1780991452919\_0020} & 5,149 & 0 & 0 & 44.9GB & 89.6GB \\
\texttt{application\_1780991452919\_0021} & 5,815 & 0 & 0 & 44.9GB & 89.4GB \\
\bottomrule
\end{tabularx}
\end{table}

从表中可以看到，正式 YARN 作业均没有 failed task 和 retried task。5\% 样本的 spill 显著高于 1\% 样本，说明本机 Docker 资源已经进入更明显的磁盘/内存交换区间；这也支持本文“不盲目承诺全量无条件通过，而用典型部分数据证明优化有效”的结论。

\subsection{典型模块 Benchmark}

根据课程实验目标与本机资源约束，后续不再启动全量 Oct+Nov 长任务，而是使用 1\% 样本中的 200,000 行典型数据补跑模块级 benchmark。该数据规模足以覆盖关联、推荐、异常和实验分流的核心计算路径，同时避免全量运行对 Docker Desktop 内存与磁盘造成不必要压力。结果如表\ref{tab:module-results}所示。

\begin{table}[htbp]
\centering
\caption{典型部分数据模块 benchmark 结果}
\label{tab:module-results}
\scriptsize
\begin{tabularx}{0.92\textwidth}{l l r r r}
\toprule
\textbf{模块} & \textbf{任务} & \textbf{输入行} & \textbf{输出行} & \textbf{耗时(s)} \\
\midrule
Affinity & \texttt{affinity\_pipeline} & 199,895 & 80 & 3.852 \\
Recommendation & \texttt{recommendation\_pipeline} & 199,895 & 1,000 & 1.740 \\
Anomaly & \texttt{anomaly\_pipeline} & 199,895 & 644 & 7.865 \\
Experimentation & \texttt{experimentation\_pipeline} & 199,895 & 21,716 & 8.703 \\
\bottomrule
\end{tabularx}
\end{table}

\subsection{结果分析}

图\ref{fig:benchmark-evidence}以生成式信息图方式概括了 benchmark 证据，精确数值以表\ref{tab:benchmark-results}为准。1\% 样本中，YARN-only CSV 用时 964.889s，明显慢于 local baseline 和后续优化组。这说明 YARN 本身提供资源调度和容错框架，但不能自动减少 CSV 扫描、shuffle、driver collect 和算法组合膨胀。启用 AQE 后耗时降至 538.657s，相对 YARN-only 约有 1.79 倍改善；加入算法限流后耗时为 535.804s，相对 YARN-only 约有 1.80 倍改善。

\begin{figure}[htbp]
    \centering
    \includegraphics[width=0.92\textwidth]{benchmark_evidence_summary.png}
    \caption{Benchmark 证据概览（由 chat-gpt-image-2 生成，精确数值见表\ref{tab:benchmark-results}）}
    \label{fig:benchmark-evidence}
\end{figure}

5\% 样本中，YARN + algorithm CSV 用时 1875.099s，YARN + Parquet 用时 1828.329s。两组都成功完成，说明 driver collect 收敛和算法上界对稳定性有效。Parquet 组在 5\% 样本中略快，说明列式输入在更大样本上开始体现优势；但由于本机 Docker 资源有限、YARN 调度开销和重复 pipeline 模块仍然存在，Parquet 收益并非线性放大。

\subsection{实验结论}

从实验结果可以得到三点结论。第一，仅将 Spark 提交到 YARN 不会自动解决性能和内存问题，YARN-only 组反而因为调度、容器和 HDFS 访问开销表现最慢。第二，AQE 和算法层限流是关键优化，能显著降低不必要 shuffle、单分区窗口和 driver 端大收集风险。第三，Parquet 输入是后续扩展到更大数据规模的必要基础，但需要与分区、列裁剪和模块级缓存一起使用。

因此，本课题最终结论应表述为：在本机 ARM64 Docker Hadoop YARN 环境中，系统已经显著降低 OOM 风险，并通过 1\%/5\% pipeline benchmark、event log 指标和 20 万行典型模块 benchmark 给出证据。按本次实验取舍，不再强行运行全量 Oct+Nov；全量处理可作为后续资源边界扩展，在更高 Docker 内存、更多 executor 或云端集群环境中继续验证。
